\section{Open questions}\label{sec:limitations}

The questions fall into two kinds.  The first three are pending further resources: the method of
Section~\ref{sec:methods} applies to them as it stands, and the runs either did not finish
within the budgets we gave them or have not yet been attempted.  They need longer runs or faster
solvers, not new ideas.  The remaining items are of other kinds: conceptual explanations we do
not have, extensions we have not attempted, a literature comparison, and further formalization.

\paragraph{Pending further resources.}
\begin{enumerate}
\item \emph{Two minima.}  Theorem~\ref{thm:conj} leaves the minimum undetermined at two
classes of special points: it is 32 or 33 at $x=\zeta/\sqrt2$ ($\zeta=e^{i\pi/3}$),
and at $x=1+i$ it lies between the general lower bound 24 and~45.  A set of fewer than 31 rays
at $x=1+i$ is not excluded.
\item \emph{Uncertified minima.}  The entries of Table~\ref{tab:catalogue} marked ``$\le$''.
\item \emph{All minimum sets on the circle.}  We exhibit one 45-word minimum set; the others have
not been enumerated.
\end{enumerate}

\paragraph{Open.}
\begin{enumerate}[resume]
\item \emph{The two circles.}  The family of 33-ray sets on $|x|^2=\tfrac12$ is that of Gould
and Aravind (Section~\ref{sec:halfcircle}).  The minima at generic points, 33 on that circle and
45 on $|x|^2=2$, are both results of the certifier.  We have no argument by hand for either, and
no explanation of why the two circles differ.  (The class $\zeta/\sqrt2$ on the inner circle
remains undetermined, item~1.)
\item \emph{Complex three-symbol alphabets.}  Theorem~\ref{thm:threesymbol} settles real
$\{0,\pm1,\pm x,\pm y\}$.  For complex $x,y$ the relations are again finitely many, but the
parameter space has four real dimensions and we have not attempted the enumeration.  Could a
complex three-symbol pool contain a KS set of fewer than 31 rays?
\item \emph{The second flex.}  Two infinitesimal flexes of the 45 words were found numerically
at both phases tested (Section~\ref{sec:rigidity}).  Is the second one a genuine deformation, leading off the
family of Theorem~\ref{thm:circle}?
\item \emph{Completion.}  Is there a finiteness criterion for completed pools?  Why does
completion make the golden-ratio and $\sqrt[3]2$ pools uncolorable but not those of $\sqrt3$ and
$\sqrt5$?
\item \emph{The gap.}  This paper does not close the gap between the lower bound 24 and
the 31 rays of Conway and Kochen.  Theorems~\ref{thm:twosymbol} and~\ref{thm:threesymbol} exclude a
smaller set with a two-symbol or a real three-symbol alphabet and nothing more; the other minima are certified within named finite pools
only.  Earlier
versions suggested that the alphabet analysis could prune the hypergraph searches
of~\cite{LiBrightGanesh2024,KirchwegerPeitlSzeider2023}; we see no basis for that, since a
hypothetical smaller set need not have a small alphabet.
\item \emph{Catalogue comparison.}  Whether the 43-, 45- and 47-ray sets appear in the automated
catalogues of three-dimensional KS sets.  (This needs reading, not computation, and has not been
done.)
\item \emph{Formalization.}  The three continuous-parameter classifications
(Theorems~\ref{thm:twosymbol}, \ref{thm:threesymbol} and~\ref{thm:conj}) are not formalized:
Lean checks Theorem~\ref{thm:6n} for $n\ge3$ and the uncolorability and minimality statements
listed in Section~\ref{sec:methods}, not the reductions to finitely many cases.  Nor does it check that each pool list is complete for its alphabet, or
part~(iii) of Theorem~\ref{thm:circle}.
\end{enumerate}

\section{Conclusion}\label{sec:conclusion}

Reading the coordinates off Peres's two cubes suggests a one-parameter question: for which $x$
does the alphabet $\{0,\pm1,\pm x\}$ generate an uncolorable pool?  The answer is now complete
(Theorem~\ref{thm:twosymbol}).  Up to the sign and inversion of $x$, which do not change the pool, it is the two classical cases
$x=\sqrt2$ and $x=2$ and nothing else for real $x$; for complex $x$ it adds only $\sqrt{-2}$ and
the primitive third and sixth roots of unity; and in no case is there a KS set of fewer than 31 rays.  One unit and one
further number along each axis cannot improve on Conway and Kochen, and neither can two further
real numbers (Theorem~\ref{thm:threesymbol}).  The more interesting
structure appears when the alphabet is enlarged.  With the conjugate
included, the uncolorable values of $x$ are no longer isolated: they are the two circles
$|x|^2=2$ and $|x|^2=\tfrac12$, four lines through the origin (without the origin), and the points
$\pm2,\pm\tfrac12$ (Theorem~\ref{thm:conj}).  On the circle $|x|^2=2$ twelve orthogonalities of the form $1+1-|x|^2=0$ make the difference; the Peres set
and its imaginary twin are the two degenerate members, up to sign, of a 45-ray family; and the 43-ray set over
$\Q(\sqrt{-7})$ sits at a point of the circle where one more relation holds.  For the
roots-of-unity alphabets the answer is exact as well: uncolorable if and only if $6\mid n$, with
the primes 2 and 3 doing different jobs.

The earlier versions of this paper claimed more than this and were wrong in ways that matter.
We have tried to make this version checkable instead: every number is produced by a named script
in exact arithmetic or is labelled numerical, every minimum is either certified by a procedure
whose soundness is argued in two paragraphs or reported as a bound, six of the minima and
Theorem~\ref{thm:6n} for $n\ge3$ are checked by the Lean kernel, and the withdrawn claims are listed in
Appendix~\ref{sec:corrections}.  Code, data and the scripts behind each table are listed in
Appendix~\ref{app:repro}; Appendix~\ref{app:produced} says how the paper was produced, including
the part played by language models.

\subsection*{Acknowledgments}

The author thanks Ad\'an Cabello for pointing out the connection between the Eisenstein alphabet
and his 33-ray set; Manuel Flores Gordillo for finding the error in the rigidity computation of
versions 1--7; Jorge Lucas Gonzalez for the supergolden construction and an exact verifier;
Markus Frembs for the trace-zero relation and for engaging with the earlier classification; and
Curtis Bright for a bibliographic correction.  The author accepts responsibility for the content,
including the errors of the earlier versions.
